\documentclass[11pt,a4paper]{article}
\usepackage[utf8]{inputenc}
\usepackage[T1]{fontenc}
\usepackage[brazil]{babel}
\usepackage{geometry}
\usepackage{enumitem}
\usepackage{xcolor}
\usepackage{titlesec}
\usepackage{fancyhdr}

\geometry{margin=2.5cm}

% Cores
\definecolor{headerblue}{RGB}{0,102,204}
\definecolor{sectiongray}{RGB}{100,100,100}

% Estilo de seções
\titleformat{\section}{\Large\bfseries\color{headerblue}}{\thesection}{1em}{}
\titleformat{\subsection}{\large\bfseries\color{sectiongray}}{\thesubsection}{1em}{}

% Cabeçalho e rodapé
\pagestyle{fancy}
\fancyhf{}
\fancyhead[L]{\textbf{SIVIRA - Enumerações}}
\fancyhead[R]{\thepage}
\renewcommand{\headrulewidth}{0.5pt}

\begin{document}

% Capa
\begin{titlepage}
    \centering
    \vspace*{3cm}

    {\Huge\bfseries Sistema Inteligente para Visualização\par}
    {\Huge\bfseries e Replanejamento de Atividades\par}
    \vspace{0.5cm}
    {\Large\bfseries (SIVIRA)\par}

    \vspace{2cm}
    {\LARGE\bfseries Documentação das Enumerações\par}

    \vspace{2cm}
    {\large Módulo: \texttt{enums/}\par}

    \vfill

    {\large \today\par}
\end{titlepage}

\tableofcontents
\newpage

% Introdução
\section{Introdução}

Este documento apresenta a documentação completa das enumerações (enums) utilizadas no sistema SIVIRA. As enumerações são organizadas em três categorias principais:

\begin{itemize}
    \item \textbf{Equipamentos}: Tipos, capacidades e características de equipamentos
    \item \textbf{Funcionários}: Funções, profissões e padrões de folga
    \item \textbf{Produção}: Itens, setores, status e políticas de produção
\end{itemize}

As enumerações garantem consistência e type-safety no sistema, definindo valores constantes utilizados em todo o código.

\newpage

% EQUIPAMENTOS
\section{Enumerações de Equipamentos}

\subsection{Capacidade Unidade (\texttt{capacidade\_unidade.py})}

Define as unidades para especificar capacidade de equipamentos.

\begin{itemize}[leftmargin=2cm]
    \item \textbf{CM2}: Área em centímetros quadrados
    \item \textbf{CM3}: Volume em centímetros cúbicos
    \item \textbf{NIVEIS}: Quantidade de níveis (prateleiras, bandejas)
    \item \textbf{OPERADORES}: Número de operadores simultâneos
    \item \textbf{GRAMAS}: Capacidade em gramas
    \item \textbf{LITROS}: Capacidade em litros
    \item \textbf{QUILOS}: Capacidade em quilogramas
\end{itemize}

\subsection{Tipo Chama (\texttt{tipo\_chama.py})}

Define os níveis de intensidade de chama para equipamentos a gás.

\begin{itemize}[leftmargin=2cm]
    \item \textbf{BAIXA}: Chama de baixa intensidade
    \item \textbf{MEDIA}: Chama de média intensidade
    \item \textbf{ALTA}: Chama de alta intensidade
\end{itemize}

\subsection{Tipo Cocção (\texttt{tipo\_coccao.py})}

Define os modos de cocção disponíveis em fornos e equipamentos de cozimento.

\begin{itemize}[leftmargin=2cm]
    \item \textbf{TURBO}: Cocção com ventilação forçada
    \item \textbf{LASTRO}: Cocção por contato direto com superfície aquecida
    \item \textbf{COMBINADO}: Cocção que combina turbo e lastro
    \item \textbf{CONTINUO}: Cocção em esteira contínua
\end{itemize}

\subsection{Tipo Embalagem (\texttt{tipo\_embalagem.py})}

Define os sistemas de embalagem disponíveis para equipamentos.

\begin{itemize}[leftmargin=2cm]
    \item \textbf{SIMPLES}: Embalagem simples sem vedação especial
    \item \textbf{VACUO}: Embalagem a vácuo para conservação
    \item \textbf{SELADORA}: Sistema de selagem térmica
\end{itemize}

\subsection{Tipo Equipamento (\texttt{tipo\_equipamento.py})}

Define os tipos de equipamentos disponíveis no sistema de produção.

\begin{itemize}[leftmargin=2cm]
    \item \textbf{MISTURADORAS}: Equipamentos para mistura de ingredientes
    \item \textbf{BANCADAS}: Superfícies de trabalho
    \item \textbf{BALANCAS}: Equipamentos de pesagem
    \item \textbf{FORNOS}: Equipamentos de cocção por calor
    \item \textbf{MISTURADORAS\_COM\_COCCAO}: Equipamentos que misturam e cozinham
    \item \textbf{FOGOES}: Equipamentos com bocas para cocção
    \item \textbf{REFRIGERACAO\_CONGELAMENTO}: Equipamentos de refrigeração
    \item \textbf{EMBALADORAS}: Equipamentos para embalagem
    \item \textbf{MODELADORAS}: Equipamentos para modelagem de massas
    \item \textbf{DIVISORAS\_BOLEADORAS}: Equipamentos para divisão e boleamento
    \item \textbf{BATEDEIRAS}: Equipamentos para bater massas
    \item \textbf{FRITADEIRAS}: Equipamentos para fritura
    \item \textbf{ARMARIOS\_PARA\_FERMENTACAO}: Equipamentos para fermentação
\end{itemize}

\subsection{Tipo Mistura (\texttt{tipo\_mistura.py})}

Define os ritmos de execução dos equipamentos de mistura.

\begin{itemize}[leftmargin=2cm]
    \item \textbf{LENTA}: Mistura em velocidade lenta
    \item \textbf{RAPIDA}: Mistura em velocidade rápida
    \item \textbf{SEMI\_RAPIDA}: Mistura em velocidade intermediária
\end{itemize}

\subsection{Tipo Pressão Chama (\texttt{tipo\_pressao\_chama.py})}

Define as configurações de pressão e chama para equipamentos a gás.

\begin{itemize}[leftmargin=2cm]
    \item \textbf{BAIXA\_PRESSAO}: Operação com baixa pressão de gás
    \item \textbf{ALTA\_PRESSAO}: Operação com alta pressão de gás
    \item \textbf{CHAMA\_UNICA}: Sistema com queimador único
    \item \textbf{CHAMA\_DUPLA}: Sistema com duplo queimador
\end{itemize}

\subsection{Tipo Velocidade (\texttt{tipo\_velocidade.py})}

Define os níveis de velocidade para equipamentos rotativos.

\begin{itemize}[leftmargin=2cm]
    \item \textbf{BAIXA}: Velocidade baixa de rotação
    \item \textbf{MEDIA}: Velocidade média de rotação
    \item \textbf{ALTA}: Velocidade alta de rotação
\end{itemize}

\newpage

% FUNCIONÁRIOS
\section{Enumerações de Funcionários}

\subsection{Tipo Folga (\texttt{tipo\_folga.py})}

Define os padrões de folga para funcionários.

\begin{itemize}[leftmargin=2cm]
    \item \textbf{DIA\_FIXO\_SEMANA}: Folga em dia fixo da semana (ex: toda segunda-feira)
    \item \textbf{DIA\_FIXO\_MES}: Folga em dia fixo do mês (ex: todo dia 15)
    \item \textbf{N\_DIA\_SEMANA\_DO\_MES}: Folga no n-ésimo dia da semana do mês (ex: segunda segunda-feira)
\end{itemize}

\subsection{Tipo Função (\texttt{tipo\_funcao.py})}

Define as funções hierárquicas disponíveis no sistema de produção.

\begin{itemize}[leftmargin=2cm]
    \item \textbf{CHEFE\_SALGADEIRO}: Responsável pelo setor de salgados
    \item \textbf{SALGADEIRO}: Profissional de produção de salgados
    \item \textbf{AJUDANTE\_DE\_COZINHA}: Assistente de cozinha
    \item \textbf{CHEFE\_CONFEITEIRO}: Responsável pelo setor de confeitaria
    \item \textbf{CONFEITEIRO}: Profissional de confeitaria
    \item \textbf{AJUDANTE\_DE\_CONFEITARIA}: Assistente de confeitaria
    \item \textbf{AJUDANTE\_DE\_SALGADOS}: Assistente do setor de salgados
\end{itemize}

\subsection{Tipo Profissional (\texttt{tipo\_profissional.py})}

Define os tipos de profissionais disponíveis no sistema.

\begin{itemize}[leftmargin=2cm]
    \item \textbf{PADEIRO}: Profissional especializado em panificação
    \item \textbf{AUXILIAR\_DE\_PADEIRO}: Assistente de padeiro
    \item \textbf{ALMOXARIFE}: Responsável pelo almoxarifado
    \item \textbf{COZINHEIRO}: Profissional de cozinha
    \item \textbf{CONFEITEIRO}: Profissional de confeitaria
    \item \textbf{AUXILIAR\_DE\_CONFEITEIRO}: Assistente de confeiteiro
\end{itemize}

\newpage

% PRODUÇÃO
\section{Enumerações de Produção}

\subsection{Dia Semana (\texttt{dia\_semana.py})}

Define os dias da semana utilizados no sistema.

\begin{itemize}[leftmargin=2cm]
    \item \textbf{SEGUNDA}: Segunda-feira
    \item \textbf{TERCA}: Terça-feira
    \item \textbf{QUARTA}: Quarta-feira
    \item \textbf{QUINTA}: Quinta-feira
    \item \textbf{SEXTA}: Sexta-feira
    \item \textbf{SABADO}: Sábado
    \item \textbf{DOMINGO}: Domingo
\end{itemize}

\subsection{Política Produção (\texttt{politica\_producao.py})}

Define as políticas de produção de itens.

\begin{itemize}[leftmargin=2cm]
    \item \textbf{SOB\_DEMANDA}: Item produzido apenas sob demanda
    \item \textbf{ESTOCADO}: Item produzido para estoque
    \item \textbf{AMBOS}: Item pode ser produzido sob demanda ou para estoque
\end{itemize}

\subsection{Status Ordem (\texttt{status\_ordem.py})}

Define os status possíveis de uma ordem de produção.

\begin{itemize}[leftmargin=2cm]
    \item \textbf{PENDENTE}: Ordem aguardando processamento
    \item \textbf{EM\_PRODUCAO}: Ordem em execução
    \item \textbf{FINALIZADO}: Ordem concluída com sucesso
    \item \textbf{CANCELADO}: Ordem cancelada
\end{itemize}

\subsection{Status Pedido (\texttt{status\_pedido.py})}

Define os status possíveis de um pedido de produção.

\begin{itemize}[leftmargin=2cm]
    \item \textbf{PENDENTE}: Pedido aguardando processamento
    \item \textbf{EM\_PRODUCAO}: Pedido em execução
    \item \textbf{FINALIZADO}: Pedido concluído com sucesso
    \item \textbf{CANCELADO}: Pedido cancelado
\end{itemize}

\subsection{Tipo Item (\texttt{tipo\_item.py})}

Define os tipos de itens na hierarquia de produção.

\begin{itemize}[leftmargin=2cm]
    \item \textbf{PRODUTO}: Item final vendável
    \item \textbf{SUBPRODUTO}: Componente intermediário de um produto
    \item \textbf{INSUMO}: Ingrediente básico
\end{itemize}

\subsection{Tipo Produto Modelado (\texttt{tipo\_produto\_modelado.py})}

Define os tipos de produtos que requerem modelagem.

\begin{itemize}[leftmargin=2cm]
    \item \textbf{PAES}: Produtos de panificação
    \item \textbf{SALGADOS}: Produtos salgados
\end{itemize}

\subsection{Tipo Setor (\texttt{tipo\_setor.py})}

Define os setores de produção no sistema.

\begin{itemize}[leftmargin=2cm]
    \item \textbf{PANIFICACAO}: Setor de panificação
    \item \textbf{SALGADOS}: Setor de salgados
    \item \textbf{CONFEITARIA}: Setor de confeitaria
    \item \textbf{ALMOXARIFADO}: Setor de armazenamento
    \item \textbf{COZINHA}: Setor de cozinha
\end{itemize}

\subsection{Unidade Medida (\texttt{unidade\_medida.py})}

Define as unidades de medida utilizadas no sistema.

\begin{itemize}[leftmargin=2cm]
    \item \textbf{GRAMAS}: Medida de massa em gramas
    \item \textbf{LITROS}: Medida de volume em litros
    \item \textbf{METROS}: Medida de comprimento em metros
    \item \textbf{CENTIMETROS}: Medida de comprimento em centímetros
    \item \textbf{MILILITROS}: Medida de volume em mililitros
    \item \textbf{UNIDADE}: Unidade contável
    \item \textbf{PORCOES}: Medida por porções
    \item \textbf{FUNCIONARIOS}: Quantidade de funcionários
    \item \textbf{UNIDADES\_POR\_HORA}: Taxa de produção
    \item \textbf{COLHERES}: Medida em colheres
    \item \textbf{XICARAS}: Medida em xícaras
\end{itemize}

\newpage

% Resumo
\section{Resumo}

\subsection{Estatísticas}

\begin{tabular}{lr}
\textbf{Categoria} & \textbf{Quantidade} \\
\hline
Enumerações de Equipamentos & 8 \\
Enumerações de Funcionários & 3 \\
Enumerações de Produção & 8 \\
\hline
\textbf{Total} & \textbf{19} \\
\end{tabular}

\vspace{1cm}

\subsection{Observações}

\begin{itemize}
    \item Todas as enumerações seguem o padrão Google Docstrings
    \item Cada enum possui documentação de módulo e de classe
    \item Os valores são type-safe e garantem consistência no sistema
    \item As enumerações são utilizadas em todo o sistema SIVIRA para padronização
\end{itemize}

\end{document}
